% Section only. Merge application-references.bib into the manuscript bibliography.
% Requires booktabs and natbib; alg:gd is defined in the main manuscript.
\section{Application Studies}\label{Sec: applications}

We study cross-hospital disease prediction, transfer of air-quality prediction
across monitoring stations, and positive-unlabeled spam ranking. In the first
two studies, client groups correspond to the hospitals and monitoring stations
identified in the original datasets. For spam, we construct a positive-unlabeled
learning task and assign emails to three clients. These studies
address two distinct questions: whether density-ratio correction improves the
downstream task, and whether distributed estimation retains the accuracy of
pooled estimation when source observations remain at their clients.

\subsection{Methods and Evaluation}

Let $P=\sum_{c=1}^C\pi_cP_c$ denote the source mixture and $Q$ the server
distribution. We compare distributed gradient descent (GD;
Algorithm~\ref{alg:gd}), closed-form random-feature KuLSIF (RF), and pooled
Gaussian-kernel KuLSIF (Kernel). Kernel uses the same fitting observations as
GD and provides a centralized accuracy reference. We also consider a local
approach (Local):
each client independently fits exact-kernel KuLSIF using only its own source
observations and the common server fitting and validation samples. No
information is exchanged among clients during local fitting. This local
baseline estimates $q/p_c$ rather than the common pooled ratio $q/p$.

All methods share the same data splits and target-only feature standardization;
missing predictors are imputed using fitting-server means. Each ratio method
selects its bandwidth and regularization separately using held-out KuLSIF
loss, without downstream evaluation outcomes. Bandwidths are
$\{0.5,1,2\}\sigma_0$, where $\sigma_0$ is the median fitting-server
pairwise distance. The hospital and air-quality studies use
$\lambda\in\{0.01,0.1\}$, $R=128$ features, and $3000$ GD updates per
candidate. The higher-dimensional spam study uses
$\lambda\in\{0.001,0.01,0.1\}$, $R=512$, and $5000$ updates. GD starts
at zero, uses step size $(2+\lambda)^{-1}$, and shares the same feature
randomness with RF.

For the two prediction studies, source observations for ratio fitting,
validation, and predictive-model training are disjoint. The server provides
covariates for ratio estimation; its held-out outcomes are used only to
evaluate the final predictions. Given $w(x)=\max\{0,\widehat r(x)\}$, the
predictive model minimizes
\[
 \frac{\sum_{c=1}^C\sum_{i=1}^{s_c}
       w(X_{ci}^{\mathrm{model}})
       \ell\{f_\theta(X_{ci}^{\mathrm{model}}),Y_{ci}^{\mathrm{model}}\}}
      {\sum_{c=1}^C\sum_{i=1}^{s_c}w(X_{ci}^{\mathrm{model}})}
 +\frac{0.01}{2}\|\theta_{-0}\|_2^2,
\]
where $\theta_{-0}$ excludes the intercept. We use linear logistic models for
disease prediction and linear regression for air quality, standardizing the
regression response on the source model-training sample. Unweighted uses
the same coordinated predictor with $w\equiv1$.

The local approach fits a separate weighted predictor at each source; its
reported loss is the source-size-weighted average of those predictors' target
losses. Consequently, its difference from GD reflects both ratio estimation
and collaboration in fitting the predictor.

We report five repetitions per hospital or station target and twenty spam
splits. Hospital scores are averaged over repetitions and then equally over
target hospitals. For air quality, we average target mean squared errors in
the same way and then take the square root. For spam, we report the mean area
under the precision--recall curve (AUPRC) across splits, computed on the combined
evaluation emails. We also report source-size-weighted within-provider AUPRC
to distinguish ranking across providers from ranking within each provider.
Repetitions reuse the same public datasets; they are not independent population
samples. The reported comparisons are descriptive rather than formal
equivalence or noninferiority tests.

\begin{table}[htbp]
\centering
\caption{Prediction at held-out hospitals and monitoring stations. Lower is
better. The local approach uses separate predictors. All four hospitals and
all twelve stations are included.}
\label{tab:applications-natural}
\medskip
\small
\setlength{\tabcolsep}{4pt}
\begin{tabular}{llrrrrr}
\toprule
Study & Metric & Unweighted & Local & GD & RF & Kernel \\
\midrule
Heart disease & Brier & 0.1730 & 0.2710 & 0.1630 & 0.1630 & 0.1676 \\
Beijing air quality & RMSE & 51.13 & 55.79 & 51.43 & 51.43 & 52.01 \\
\bottomrule
\end{tabular}
\end{table}

\subsection{Cross-Hospital Heart Disease Prediction}

The Heart Disease data contain records from Cleveland, Hungary, Switzerland,
and VA Long Beach~\citep{janosi1989heart}. We use all four processed files,
comprising $920$ patient records, and define disease presence as a recorded
diagnosis greater than zero. Each hospital serves as the target in turn,
with the other three hospitals as source clients. Predictors are age, sex,
resting blood pressure, cholesterol, maximum heart rate, exercise-induced
angina, ST depression, and fixed chest-pain indicators. Unknown values and
zero blood-pressure, cholesterol, or heart-rate values are treated as missing.

Patients are randomly assigned within their original hospital to model
training ($25\%$), ratio fitting ($25\%$), ratio validation ($15\%$), and
evaluation ($35\%$). For computational screening, each source contributes
at most $80$ ratio-fitting, $80$ validation, and $160$ model-training records;
the target contributes at most $240$ fitting, $80$ validation, and $480$
evaluation records. The available hospital samples are often smaller than
these caps. No target outcome is used to train the predictor.

Table~\ref{tab:applications-natural} shows a mean target Brier score of
$0.1630$ for GD, compared with $0.1730$ for Unweighted, $0.2710$ for the local
approach, and $0.1676$ for Kernel. The improvement over Unweighted is approximately
$5.8\%$. This average masks variation across hospitals: for Switzerland,
GD reduces Brier score from $0.2237$ to $0.1558$, whereas Unweighted performs
better at the other three targets. Thus the average benefit is concentrated
in one hospital. Hospital-specific missingness and differences in the conditional disease
distribution remain possible, so the study does not verify a pure
covariate-shift model or establish clinical deployment performance.

\subsection{Air-Quality Prediction Across Monitoring Stations}

The Beijing Multi-Site Air Quality data contain hourly observations from
twelve monitoring stations between March 2013 and February
2017~\citep{chen2017beijing}. We predict measured daily mean PM$_{2.5}$ from
temperature, pressure, dew point, precipitation, wind speed, and sine/cosine
seasonal terms. Daily means require at least eighteen observed hours per
variable. Days without an observed response are excluded, yielding $17164$
station-days; missing predictors are imputed as described above. Other
pollutants and station identifiers are excluded from the predictors.

Each station is held out in turn, leaving eleven source clients. Calendar
months are ordered into model training (first $25\%$), ratio fitting (next
$35\%$), validation (next $10\%$), and evaluation (last $30\%$), with one
complete month removed at each boundary. These roles are shared across
stations. Consequently, the task involves both station transfer and temporal
change. We use the same sample caps as in the hospital study, sampling within
the assigned roles. Temporal separation avoids reusing neighboring days
across role boundaries but does not establish independence.

GD attains an aggregate target RMSE of $51.43\,\mu\mathrm{g}/\mathrm{m}^3$,
compared with $52.01$ for Kernel and $55.79$ for the local approach.
Its accuracy matches RF to the displayed precision.
However, Unweighted achieves the lowest aggregate RMSE, $51.13$.
This application therefore shows that distributed ratio estimation retains
accuracy close to the pooled references at held-out stations, without an
improvement from importance weighting. Spatial
dependence, temporal change, and omitted pollution drivers can limit the
usefulness of a correction based only on the observed covariates.

\subsection{Positive-Unlabeled Spam Ranking}

We use Spambase, which contains $4601$ labeled emails represented by $57$
word-frequency, character-frequency, and capitalization
features~\citep{hopkins1999spambase}. We construct a setting in which the
server holds verified spam and three clients hold unlabeled emails.
Training uses no labeled non-spam examples. Labels identify the server's
positive reference sample and are otherwise reserved for evaluation.
The email features and labels are observed, but provider allocation and
the positive-unlabeled observation protocol are constructed.

Let $q_+(x)$ be the positive-class density, $p(x)=\sum_c\pi_cp_c(x)$ the
pooled email density, and $\alpha=\Pr(Y=1)$ the pooled spam prevalence.
If verified positives represent the pooled positive class, Bayes' rule gives
\[
 \Pr(Y=1\mid X=x)=\alpha\,\frac{q_+(x)}{p(x)}.
\]
Thus estimating $q_+/p$ ranks spam without knowing $\alpha$. We make no
claim of calibrated posterior probabilities. Local models instead estimate
$q_+/p_c$, whose scores need not share the pooled ranking across providers.

For each of twenty splits, disjoint pools supply up to $600$ unlabeled
source records, $600$ verified positives, and $1200$ evaluation emails.
The source and positive pools are each split $80\%/20\%$ for fitting and
validation; the positive pool is typically smaller than its cap.
Duplicate-record groups remain in a single pool. Heterogeneous provider
ownership is assigned probabilistically using capitalization style, without
labels, with at least $3\%$ assignment probability to every provider for
every covariate value. A random-ownership control reallocates the same
fitting, validation, and evaluation records while fixing the pooled problem
and feature randomness. Local models score only their own provider's
evaluation records. A positive-only Gaussian density estimator (KDE),
selected by held-out positive log likelihood, assesses whether estimating
the email-density denominator improves ranking. Scores are untruncated.

\begin{table}[htbp]
\centering
\caption{Spam ranking: mean AUPRC over twenty splits; higher is better.
KDE uses only the positive reference density. Changing ownership preserves
the pooled task, so GD, RF, Kernel, and KDE have identical global results.}
\label{tab:applications-spam}
\medskip
\small
\begin{tabular}{lrrrrr}
\toprule
Provider allocation & KDE & Local & GD & RF & Kernel \\
\midrule
Random & 0.6656 & 0.8338 & 0.8619 & 0.8611 & 0.8730 \\
Heterogeneous & 0.6656 & 0.7720 & 0.8619 & 0.8611 & 0.8730 \\
\bottomrule
\end{tabular}
\end{table}

Under heterogeneous ownership, GD attains AUPRC $0.8619$, compared with
$0.6656$ for KDE and $0.7720$ for the local approach
(Table~\ref{tab:applications-spam}). Its gap from Kernel is $0.0111$.
Reallocating the same data randomly raises the local approach's AUPRC to $0.8338$ while
leaving the global estimators unchanged. Within-provider AUPRC under
heterogeneous ownership is $0.8476$ for GD and $0.8394$ for the local approach; this
smaller difference indicates that much of the global gain concerns score
comparability across providers. The study illustrates the value of a common
pooled denominator for global ranking, under the assumption of representative
verified positives.

\begin{samepage}
Across the three studies, GD retains accuracy close to pooled KuLSIF while
keeping source observations local. Weighting improves average hospital
prediction and global spam ranking, but not air-quality prediction.
\end{samepage}
